\section{Callais Compliance Architecture}
\label{sec:callais}

\subsection{The Disentanglement Requirement}

\callais{}, 608 U.S.\ \_\_\_ (2026), at p.~36, articulates the disentanglement
requirement as follows:
a redistricting process that combines race-conscious and partisan signals in
the same objective function is constitutionally suspect because ``a court
cannot isolate the racial motivation from the partisan motivation in evaluating
the resulting map's constitutionality.''
The majority did not hold that either signal is impermissible; it held that
the two signals cannot be \emph{combined} in a way that makes their individual
contributions undecipherable.

This requirement maps precisely onto the toolbox architecture.
The toolbox has exactly two signal-bearing modes:
\begin{itemize}
  \item \textbf{VRASection (B.14)}: Geographic alignment score on minority VAP.
    Permissible under \callais{} because (a) it addresses a geographic property
    (where minority population is concentrated), (b) it enters only as a ratio
    selector, not as a METIS constraint, and (c) compactness is the predominant
    criterion at $w_\text{vra} = 0.40$.
  \item \textbf{ProportionalSection (B.12, planned)}: Partisan vote-share balance
    constraint.
    Permissible under \rucho{}'s state-court exception because it is a partisan
    signal, not a racial signal, and it addresses the proportionality criterion
    directly.
    Mutex-gated against VRASection.
\end{itemize}

\subsection{How VRASection Avoids Racial Predominance}

VRASection's design makes three specific architectural choices to stay below
the \shaw{}--\miller{} racial-predominance threshold.

\paragraph{Choice 1: Geographic proxy, not racial count.}
The alignment score $A = |\text{MVAP}_\text{frac}(L) - \text{MVAP}_\text{frac}(R)|$
measures how much minority VAP is concentrated on one side of a proposed
bisection.
This is a \emph{geographic property} of the existing population distribution,
not a racial target.
The algorithm does not specify ``the left side must contain at least $X\,\%$
minority population''; it prefers bisections that \emph{happen to place} a
compact minority community on one side because that is what minimises the
score function.
The geographic structure drives the decision; the racial measurement merely
\emph{detects} that structure.

\paragraph{Choice 2: Compactness-dominated weight.}
At $w_\text{vra} = 0.40$, the score function is:
\[
  \text{score}(\text{ratio}) = \underbrace{\mathrm{EC}_\text{norm}}_{\text{compactness}} -
  \underbrace{0.40 \cdot A \cdot \max(\mathrm{EC}_\text{norm}, 1)}_{\text{alignment bonus}}
\]
The alignment bonus can reduce the compactness penalty by at most $0.40 \times A$,
where $A \leq 1$.
The minimum alignment effect is zero; the maximum is a $40\,\%$ reduction in
the compactness score.
In practice, for Alabama ($A = 0.42$ at the $2:5$ winning ratio), the alignment
bonus reduces the score from 448.7 to 322.5---a $28\,\%$ reduction.
The compactness penalty (448.7) is still the dominant term; the alignment bonus
is a correction, not a driver.

This weight satisfies the \miller{} test as interpreted in \case{Bush v.\ Vera}
(1996): ``race was not the predominant factor overriding traditional districting
principles.''
Traditional compactness (as measured by normalised edge-cut) is the dominant
criterion at $w_\text{vra} = 0.40$.

\paragraph{Choice 3: No MetisVra failure mode.}
The MetisVra approach \citep{b3paper} added minority VAP as a second METIS
balance constraint.
This fails in 3 of 5 majority-minority states due to numerical scale mismatch:
minority VAP fractions (0.0--1.0) and census population counts (10,000--200,000)
differ by $60$--$800\times$.
More fundamentally, MetisVra \emph{racially targets}: it specifies an explicit
minority fraction that the district must achieve.
VRASection avoids this by never specifying a target minority fraction.
The alignment score measures geographic concentration; it does not mandate a
racial outcome.

\subsection{The CLI Mutex Gate}

The \texttt{callais\_preflight} function is the runtime enforcement mechanism.
It is called at the start of every \texttt{redist state} and \texttt{redist run}
invocation, before any data is loaded:

\begin{verbatim}
fn callais_preflight(args: &StateArgs) -> Result<()> {
    if args.partition_mode == VraSection
       && args.proportional_section {
        return Err(CallaísViolation {
            message: "VRASection and ProportionalSection cannot be \
                      combined in one run (Callais p.36). Run each \
                      mode independently.",
            citation: "Louisiana v. Callais, 608 U.S. ___ (2026), p.36",
        });
    }
    Ok(())
}
\end{verbatim}

The function also checks for weaker disentanglement violations: if the user
sets $w_\text{vra} > 0$ while also specifying \texttt{--election-year}
(which enables election-result analysis), a warning is emitted:

\begin{verbatim}
  [BOUNDARY] callais_warning: VRASection with --election-year set.
  Election results will be recorded in the manifest but were NOT
  used as input to the partition algorithm. Confirm this is correct.
\end{verbatim}

This warning is not a blocking error---election results can legitimately be
recorded for post-hoc analysis---but it documents that the partisan data was
not part of the optimisation.

\subsection{The Plan Manifest as Audit Record}

Every \texttt{redist} run produces a \texttt{whatif-manifest v1} JSON sidecar
with the following fields relevant to Callais compliance:

\begin{verbatim}
{
  "partition_mode": "vra-section",
  "w_vra": 0.40,
  "proportional_section": false,
  "election_year": null,
  "callais_preflight_passed": true,
  "callais_preflight_version": "2026-04-30",
  "inputs": {
    "adj_graph_sha256": "a3f7...",
    "minority_vap_source": "Census P4 table, 2020",
    "partisan_data_source": null
  },
  "overrides_hash": "d5e2...",
  "log_sha256_chain": ["...", "..."]
}
\end{verbatim}

A reviewing court or opposing expert can verify:
(a) the partition mode is VRASection but not ProportionalSection,
(b) partisan data was not an input (\texttt{partisan\_data\_source: null}),
(c) the \texttt{callais\_preflight} check was passed (not bypassed),
(d) the adjacency graph SHA-256 matches the publicly released data,
(e) the overrides hash is consistent with the stated parameters, and
(f) the log SHA-256 chain has not been tampered with (single-byte tamper
    is detectable).

\subsection{The Strong-Inference Test Applied to the Toolbox}

\callais{}'s ``strong-inference'' framework (holding 2) requires \S2 challengers
to demonstrate that the state's stated political goals could have been achieved
\emph{with better minority outcomes}.
The toolbox directly provides this counterfactual analysis.

\paragraph{Procedure.}
\begin{enumerate}
  \item Run GeoSection (no minority signal) on the disputed state.
    Record D seat count and MM district count. This is the ``pure political
    goals'' baseline.
  \item Run VRASection ($w_\text{vra} = 0.40$) on the same state.
    Record D seat count and MM district count.
  \item If VRASection achieves at least as many D seats as GeoSection \emph{and}
    more MM districts, the strong-inference test is satisfied: the state's
    political goals could have been achieved with better minority outcomes.
\end{enumerate}

For Alabama:
\begin{itemize}
  \item GeoSection (stated political goal: efficient district layout):
    0D, 1 MM district.
  \item VRASection: 1D, \textbf{2 MM districts} (confirmed, B.14).
  \item Strong-inference test: \textbf{satisfied}.
    The state's goals (compact, population-balanced districts) could have been
    achieved with two minority-opportunity districts.
    The \allenmilligan{} remedy is computable from the toolbox.
\end{itemize}

For Georgia:
\begin{itemize}
  \item GeoSection: 5D, 1 MM district (estimated).
  \item VRASection: 5D, 3 MM districts (estimated).
  \item Strong-inference test: satisfied (same seat count, better minority outcomes).
\end{itemize}

For North Carolina:
\begin{itemize}
  \item GeoSection: 5D, 1 MM district.
  \item VRASection: 5D, 2 MM districts (estimated, Charlotte/Raleigh isolation).
  \item Strong-inference test: satisfied.
\end{itemize}

\subsection{Separation of Legal Signals: A Formal Statement}

\begin{proposition}[Callais compliance of the toolbox]
\label{prop:callais}
The algorithmic redistricting toolbox satisfies the \callais{} p.~36
disentanglement requirement if and only if the following conditions hold:
\begin{enumerate}
  \item At most one of $\{$VRASection, ProportionalSection$\}$ is active in
    any single \texttt{redist} run.
  \item The \texttt{callais\_preflight} gate is not bypassed (no
    \texttt{--skip-preflight} flag).
  \item The plan manifest \texttt{whatif-manifest v1} is produced and its
    hash chain is unbroken.
\end{enumerate}
Under these conditions, a reviewing court can determine from the manifest alone
which signal was used, verify that no other signal was present as an input,
and reproduce the run from the stated parameters.
\end{proposition}

\begin{proof}
Conditions 1--2 prevent the simultaneous use of race-conscious and partisan
signals.
The \texttt{callais\_preflight} function is called before any data is loaded;
it cannot be bypassed without modifying the binary.
Condition~3 provides the audit record: the SHA-256 hash of the adjacency graph
verifies the data provenance; the \texttt{overrides\_hash} verifies the
parameters; the log hash chain verifies post-hoc integrity.
A court can reproduce the run by providing the adjacency graph (hash-verified)
and the stated parameters, and checking that the resulting partition matches
the produced map.
\end{proof}
